% !TEX ROOT = ../Thesis.tex 
%-
%-> 中文摘要
%-
\chapter{摘\quad 要}\chaptermark{摘\quad 要}
\linespread{1.5}
\zihao{-4}

在大语言模型与检索增强生成（Retrieval-Augmented Generation, RAG）技术迅猛发展的背景下，多轮对话检索已成为对话式人工智能系统落地应用的关键环节。然而，真实多轮对话中普遍存在的指代、省略、话题漂移等语言现象，使得当前轮次的查询往往语义残缺，无法被现有检索器直接理解和召回。查询重写（Query Rewriting）作为在检索之前将上下文依赖的用户问句改写为语义自包含独立查询的关键预处理模块，直接决定了下游检索与生成的上限。然而，现有查询重写方法仍面临三大挑战：其一，对长历史、跨话题对话中的噪声轮次缺乏显式建模，导致重写时引入大量无关信息；其二，指代消解与省略补全被混合为端到端序列生成任务，训练目标不明确且错误难以定位；其三，重写目标局限于"像人写的参考重写"，与下游检索器的真实偏好存在系统性偏差。针对上述痛点，本文系统研究了面向多轮对话检索的查询重写方法，提出从"历史剪枝—渐进式重写—检索反馈优化"三阶段构建查询重写框架，并在多个公开基准上进行了深入的实证分析。本文的主要研究工作与创新成果如下：

首先，针对多轮对话中历史轮次的话题漂移问题，提出了一种话题漂移感知的对话历史剪枝方法。不同于将完整历史一视同仁地拼接输入的传统做法，本方法在轮次级引入三分类相关性判别器（强相关 / 弱相关 / 无关），基于预训练语言模型构建轻量级分类头，并利用公开对话数据集中话题切换标注进行弱监督训练。该模块将输入给重写器的上下文压缩为高相关子序列，显著降低输入长度与噪声密度。在 TopiOCQA 与 QReCC 的验证实验表明，剪枝后的历史使下游重写器的 BLEU 指标提升 3.6 个百分点、Recall@20 提升 4.1 个百分点。

其次，针对指代消解与省略补全耦合导致错误难以分离的问题，提出了一种两阶段渐进式查询重写框架（Progressive Rewriter, PRW）。该框架将重写解耦为"指代消解"与"省略补全"两个串行子任务：第一阶段仅将指代词替换为对应实体，保持表层结构不变；第二阶段在此基础上补全缺失的主语、谓语与修饰成分，生成自包含查询。两阶段分别使用合成监督数据进行有监督微调，并通过阶段一致性约束防止语义漂移。实验证明，两阶段解耦显著提升了重写的可解释性与稳定性，相较于端到端基线在四个对话检索基准上平均带来 2.8 个百分点的 Recall@5 增益。

最后，针对重写目标与下游检索器偏好失配的问题，提出了检索反馈驱动的强化微调方法（Retrieval-Feedback Rewriter, RFR）。本文将检索器对重写查询的 Recall@k 作为奖励信号，构建"胜出重写-落败重写"的偏好对，采用直接偏好优化（Direct Preference Optimization, DPO）对重写器进行轻量化微调。该方法无需额外标注，训练稳定且成本低廉。在 QReCC、QuAC、Doc2Dial、TopiOCQA 四个公开基准上，RFR 使 Qwen2.5-3B 参数量级的重写器达到 72B 规模 zero-shot 重写器的检索性能，同时推理延迟降低一个数量级，体现出显著的实用价值。

综上所述，本文围绕多轮对话检索中的查询重写问题，构建了一套"噪声过滤—结构化重写—下游反馈"相互衔接的技术方案，深入研究了对话历史建模、渐进式生成与基于检索反馈的偏好优化等核心问题，为工业级对话式 RAG 系统提供了一种轻量、高效且鲁棒的查询重写范式。

\keywords{多轮对话检索；查询重写；检索增强生成；直接偏好优化；大语言模型}

%-
%-> 英文摘要
%-
\chapter{Abstract}\chaptermark{Abstract}

Against the backdrop of the rapid progress of large language models (LLMs) and retrieval-augmented generation (RAG), multi-turn conversational retrieval has become a critical building block for deploying conversational AI systems. In realistic multi-turn dialogues, however, linguistic phenomena such as coreference, ellipsis, and topic drift are ubiquitous, rendering the current-turn query semantically incomplete and hard for existing retrievers to understand. Query rewriting (QR), which transforms a context-dependent utterance into a self-contained standalone query before retrieval, directly determines the upper bound of downstream retrieval and generation. Nevertheless, existing QR approaches still suffer from three key limitations: (i) they lack explicit modeling of noisy turns in long, topic-switching histories, so that irrelevant information is easily leaked into the rewrite; (ii) coreference resolution and ellipsis completion are conflated into a single end-to-end sequence generation task with ill-defined training signals; and (iii) the training objective is tied to "human-written references" and is systematically misaligned with the true preferences of downstream retrievers. To address these issues, this thesis systematically investigates query rewriting methods for multi-turn conversational retrieval, and proposes a three-stage framework consisting of \emph{history pruning}, \emph{progressive rewriting}, and \emph{retrieval-feedback optimization}. The main contributions are as follows.

First, to tackle topic-drift in long conversational histories, we propose a \textbf{topic-drift-aware history pruning} module. Instead of concatenating the whole history, the method introduces a turn-level three-way relevance classifier (strongly-, weakly-, and un-related) built on top of a lightweight pretrained encoder, trained with weak supervision derived from topic-switch labels in public dialogue corpora. The pruned context is passed to the rewriter, significantly reducing input length and noise density. Experiments on TopiOCQA and QReCC show that the pruning module improves the rewriter's BLEU by 3.6 points and downstream Recall@20 by 4.1 points.

Second, to resolve the coupling between coreference and ellipsis phenomena, we propose a \textbf{Progressive Rewriter (PRW)} framework that decomposes rewriting into two sequential subtasks: a coreference-only stage that replaces anaphoric mentions while preserving surface structure, and an ellipsis-completion stage that restores missing subjects, predicates, and modifiers. The two stages are trained via supervised fine-tuning with stage-consistency constraints. The decoupled design yields better interpretability and stability, delivering an average 2.8-point Recall@5 gain over strong end-to-end baselines across four conversational retrieval benchmarks.

Third, to bridge the gap between human-style rewrites and retriever-friendly rewrites, we design a \textbf{Retrieval-Feedback Rewriter (RFR)} that treats the retriever's Recall@k as a reward signal. Winning/losing rewrite pairs are collected to form preference data, and Direct Preference Optimization (DPO) is applied for lightweight alignment of the rewriter. No additional human annotations are needed. On four public benchmarks (QReCC, QuAC, Doc2Dial, TopiOCQA), RFR enables a 3B rewriter to match or surpass a 72B zero-shot rewriter in retrieval quality while cutting inference latency by an order of magnitude.

In summary, this thesis presents a unified technical route that couples noise filtering, structured rewriting, and downstream feedback, and provides a lightweight, efficient and robust query-rewriting paradigm for industrial-scale conversational RAG systems.

\englishkeywords{multi-turn conversational retrieval; query rewriting; retrieval-augmented generation; direct preference optimization; large language models}
